\documentclass[10pt,a4paper]{amsart}
\usepackage[margin=19mm]{geometry}
\usepackage{amsmath,amssymb,mathtools}
\usepackage[scr=rsfs,cal=euler]{mathalfa}
\usepackage{xfrac}
\usepackage[unicode,hidelinks,bookmarks=false]{hyperref}
\newcommand{\ie}{i.e.\ }
\newcommand{\cf}{cf.\ }
\newcommand{\resp}{respectively\ }
\newcommand{\Subsection}{\subsection}
\providecommand{\nobreakdash}{\nobreak\hbox{-}}
\setlength{\emergencystretch}{2em}
\multlinegap0pt
\usepackage{siunitx}
\usepackage{siunitx}
\usepackage{siunitx}
\usepackage{siunitx}
\usepackage{amssymb}
\usepackage{enumerate}
\usepackage{xcolor}
\usepackage{float}
\usepackage{bm}
\usepackage{tikz}
\usepackage{algorithm}\usepackage{algpseudocode}
\usepackage{mathtools}
\usepackage{siunitx}
\newtheorem{proposition}{Proposition}
\def\mA{{\bm{A}}}
\def\mB{{\bm{B}}}
\def\mM{{\bm{M}}}
\def\mN{{\bm{N}}}
\def\va{{\vec{a}}}
\def\vb{{\vec{b}}}
\def\vc{{\bm{c}}}
\def\vd{{\bm{d}}}
\renewcommand{\vec}[1]{\mathbf{#1}}
\def\vh{{\vec{h}}}
\def\vk{{\bm{k}}}
\def\vx{{\vec{x}}}
\title{A Computational Tropical Geometry Framework for Neural Networks}
\author{Paul Lezeau, Thomas Walker, Yueqi Cao, Shiv Bhatia, Anthea Monod}
\date{Source: arXiv:2405.20174v3 (CC BY 4.0)}
\begin{document}
\maketitle

\noindent\emph{Source and licence.} A Computational Tropical Geometry Framework for Neural Networks. Version arXiv:2405.20174v3 (CC BY 4.0), distributed by arXiv under the Creative Commons Attribution 4.0 International licence, http://creativecommons.org/licenses/by/4.0/, as recorded in the arXiv licence metadata for this exact version and shown on the abstract page https://arxiv.org/abs/2405.20174v3. Full licence notice: \texttt{licenses/arxiv-2405.20174-CC-BY-4.0.txt}.\par\smallskip

\noindent\textit{Editorial scope.} Counted unit: the article's proposition eq:rf-hpq (source0.tex lines 1192--1198). The article's complete original proof of it is delivered with it (source0.tex lines 1200--1241, one \texttt{proof} environment of 42 lines). No mathematics is written, repaired or re-derived: the statement and the proof are the article's own lines, in the article's own notation, and no retained line is substituted or re-flowed.  No further source-local context is needed: every cross-reference of every retained line resolves inside the two delivered spans.  Nothing was omitted from any retained span, so the package carries no omission record.  No retained line contains a citation, so the wrapper carries no bibliography and no bibliography entry.  No retained line contains a figure environment, an image-inclusion command or drawing code, so no image is delivered and the package declares no asset dependency.  Editorial formatting only, in addition: an AMS \texttt{amsart} preamble that restates the article's own declarations and macros used by the retained lines, namely the environments \texttt{proposition} on separate counters, together with the standard packages those lines need.  The retained environments print in this document as Proposition 1 in order of appearance, whereas the article prints the same items with the numbering of its own full document.  The complete native source of the article as distributed in the arXiv e-print for this exact version is bundled unchanged as \texttt{sources/160033/source0.tex}.
\begin{proposition}
\label[proposition]{prop:effective-radius}
    Let $f = p\oslash q$ be a rational signomial. For any $\vx\in\mathbb{R}^n$, we have
    \begin{equation}\label{eq:rf-hpq}
       R_f(\vx)\le H(p,q)\max\{p(\vx)-\check{p}(\vx),\, q(\vx)-\check{q}(\vx)\}.
    \end{equation}
\end{proposition}

\begin{proof}
The polyhedra defined by 
$$
\begin{pmatrix}
            \mM_i \\
            \mN_j
        \end{pmatrix} \vx \le \begin{pmatrix}
            \vh_i \\
            \vk_j
        \end{pmatrix}
$$
of dimension $n$ form a convex refinement of linear regions of $f$. Let 
$$\mathrm{res}_{i,j}(\vx):=\begin{pmatrix}
            \mM_i \\
            \mN_j
        \end{pmatrix} \vx- \begin{pmatrix}
            \vh_i \\
            \vk_j
        \end{pmatrix}
$$        
denote the residual of $\vx$ to the polyhedron. We have 
$$
R_f(\vx)\le H(p,q)\max\left\{\|\mathrm{res}_{i,j}(\vx)_+\|_\infty: 1 \le i \le k, 1 \le j \le \ell, \dim U_i \cap V_j = n \right\}.
$$
Note that
\begin{align*}
    \|\mathrm{res}_{i,j}(\vx)_+\|_\infty
        &= \left\|\left(
            \begin{bmatrix}\mA \vx+\vc-\mathbf{1}(\langle \va_i, \vx \rangle+c_i)\\
                            \mB \vx+\vd-\mathbf{1}(\langle \vb_j, \vx \rangle+d_j)\end{bmatrix}\right)_+\right\|_\infty\\
        &=\max_{r,s}\left\{(\mA \vx+\vc)_r-(\langle \va_i, \vx \rangle+c_i),\,(\mB \vx+\vd)_s-(\langle \vb_j, \vx \rangle+d_j),\, 0\right\}\\
        &=\max\left\{p(\vx)-(\langle \va_i, \vx\rangle+c_i),\, q(\vx)-(\langle \vb_j, \vx\rangle + d_j),\, 0\right\}
\end{align*}
Therefore,
\begin{align*}
    \max_{\substack{i, j \text{ s.t} \\ \dim (U_i \cap V_j) = n}}\|\mathrm{res}_{i,j}(\vx)_+\|_\infty
        &=\max_{\substack{i, j \text{ s.t} \\ \dim (U_i \cap V_j) = n}}\left\{p(\vx)-(\langle \va_i, \vx\rangle +c_i),\, q(\vx)-(\langle \vb_j,  \vx \rangle+d_j),\, 0\right\}\\
        &\le\max\left\{p(\vx)-\min_i\{\langle \va_i, \vx \rangle +c_i\},\, q(\vx)-\min_j\{\langle \vb_j, \vx \rangle +d_j\},\, 0\right\}\\
        &=\max\left\{p(\vx)-\check{p}(\vx),\,q(\vx)-\check{q}(\vx)\right\}
\end{align*}
which proves \eqref{eq:rf-hpq}.
\end{proof}

\end{document}
